\documentclass[12pt]{article}
\usepackage[utf8]{inputenc}
\usepackage[T1]{fontenc}
\usepackage[margin=1in]{geometry}
\usepackage{setspace}
\usepackage{xcolor}
\usepackage[colorlinks=true, urlcolor=blue!60!black, linkcolor=blue!60!black]{hyperref}
\setlength{\parindent}{2em}
\setlength{\parskip}{1.25em}
\renewcommand{\baselinestretch}{1.0}
\begin{document}
\begin{center}
{\LARGE \textit{Smith v. Doe}, 538 U.S. 84 (2003)\par}
\vspace{0.8em}
{\large Case Brief --- Four-Page Edition\par}
\end{center}
\vspace{1em}

\section{Memory Jogger}

Alaska applied its sex offender registration and Internet notification law to offenders convicted before it was enacted; the Court held 6--3 that the law was a civil, nonpunitive regulation under the two-step intent--effects test, so its retroactive application did not violate the Ex Post Facto Clause (\href{https://www.law.cornell.edu/supct/html/01-729.ZO.html}{Cornell LII}).

\section{Detailed Case Facts}

In 1994 Alaska enacted the Alaska Sex Offender Registration Act (ASORA), a ``Megan's Law'' of the kind Congress encouraged through the Jacob Wetterling Act. It applied retroactively to any sex offender or child kidnapper in Alaska. \textbf{Registration:} offenders had to give police their name, aliases, address, employer, vehicle, photograph, fingerprints, and treatment history; single nonaggravated offenders verified annually for 15 years, and aggravated or repeat offenders quarterly for life. Knowing noncompliance was a crime. \textbf{Notification:} the Department of Public Safety kept fingerprints and treatment records confidential but published names, photographs, home and work addresses, vehicles, convictions, and compliance status on a public website.

John Doe I and John Doe II pleaded \textit{nolo contendere} to aggravated sexual abuse of a minor years before ASORA, were released in 1990, and completed rehabilitation programs.

\section{Procedural History}

With Jane Doe, they sued Alaska officials under 42 U.S.C. \S~1983. The district court granted summary judgment for the State. The Ninth Circuit reversed, agreeing that the legislature intended a civil scheme but holding that its punitive effects overrode that intent. The Supreme Court reversed and remanded on March 5, 2003.

\section{Judicial Votes}

Reversed and remanded, \textbf{6--3}.
\begin{itemize}
\item \textbf{Majority:} Justice Kennedy, joined by Chief Justice Rehnquist and Justices O'Connor, Scalia, and Thomas.
\item \textbf{Concurrences:} Justice Thomas; Justice Souter (in the judgment).
\item \textbf{Dissents:} Justice Stevens; Justice Ginsburg, joined by Justice Breyer.
\end{itemize}

\section{Holding}

Because the Alaska Legislature intended a civil, nonpunitive regime and respondents did not show by the ``clearest proof'' that its effects negated that intent, retroactive application of ASORA does not violate the Ex Post Facto Clause.

\section{Analysis of Opinions}

\textbf{Majority (Kennedy): intent.} Under \textit{Kansas v. Hendricks} and \textit{United States v. Ward}, and the premise of \textit{Calder v. Bull} that the Clause bars only retroactive criminal punishment, the first question is statutory construction. ASORA named ``protecting the public from sex offenders'' as its ``primary governmental interest.'' Its placement partly in the criminal procedure code and the plea-colloquy notices were not dispositive, since the notification provisions sat in the Health, Safety, and Housing Code and an administrative agency ran the program.

\textbf{Effects.} Of the \textit{Kennedy v. Mendoza-Martinez} factors: registries are recent and unlike colonial shaming or banishment, and any stigma is a collateral consequence of conviction; ASORA imposes no physical restraint, and registrants may change jobs or homes freely; incidental deterrence does not make a law punitive, and its tiers track recidivism risk; its rational connection to public safety, the ``most significant'' factor, favored Alaska; and, citing \textit{McKune v. Lile} on recidivism risk that is ``frightening and high,'' the State may act categorically, without individualized hearings.

\textbf{Thomas (concurring).} Only ``the statute on its face'' matters, not Internet implementation.

\textbf{Souter (concurring in the judgment).} With no civil label and split codification, the indicators were in ``rough equipoise''; he concurred only on the presumption of constitutionality.

\textbf{Stevens (dissenting).} A sanction imposed on everyone convicted of a crime, and no one else, that severely impairs liberty is punishment.

\textbf{Ginsburg (dissenting).} ASORA resembles supervised release and shaming and is excessive, with lifetime quarterly reporting and no path off the registry for rehabilitated offenders.

\textbf{Later development.} \textit{Ellingburg v. United States} (2026) applied \textit{Smith}'s first step to hold MVRA restitution criminal punishment.

\section{Examples: Future Applications}

\begin{enumerate}
\item \textbf{Same-Side (Civil): A Fraud Registry.} A civil securities statute publishes convicted fraudsters' photos and addresses online and requires lifetime account reporting. \textit{Smith}'s ``clearest proof'' standard heavily favors the government.
\item \textbf{Same-Side (Civil): Civil Commitment.} A probate-code commitment law for sexually violent predators, with release when no longer dangerous, follows \textit{Hendricks}.
\item \textbf{Opposite-Side (Punishment): Court-Enforced Residency Bans.} Retroactive exclusion zones imposed by a sentencing court with re-incarceration for violations reveal punitive intent under \textit{Ellingburg}.
\item \textbf{Opposite-Side (Punishment): Harsher Guidelines.} Sentencing under later, harsher Guidelines violates the Clause under \textit{Peugh v. United States} (2013).
\item \textbf{Fence-Sitter: Lifetime GPS Monitoring.} A civil-code law requires corrections-supervised GPS monitors for life, resembling the supervision \textit{Smith} distinguished.
\end{enumerate}

\section{Critique}

\textbf{The ``clearest proof'' trap.} Echoing Souter and Ginsburg, critics argue that once a legislature invokes ``public safety,'' Step Two becomes nearly insurmountable, inviting harsh collateral penalties relabeled as ``civil.''

\textbf{Empirical fragility and a dated analogy.} The ``frightening and high'' recidivism claim traces to a non-peer-reviewed \textit{Psychology Today} article later disavowed by its author. And treating a searchable, indexed registry as ``a visit to an official archive'' understates permanent digital stigma and private ostracism.

\section{Key Quotations}

\begin{quote}
``If the intention of the legislature was to impose punishment, that ends the inquiry. If, however, the intention was to enact a regulatory scheme that is civil and nonpunitive, we must further examine whether the statutory scheme is `so punitive either in purpose or effect as to negate [the State's] intention' to deem it `civil.'\,''
\end{quote}
\noindent --- Justice Kennedy, for the Court, 538 U.S. at 92.

\begin{quote}
``The excessiveness inquiry of our ex post facto jurisprudence is not an exercise in determining whether the legislature has made the best choice possible to address the problem it seeks to remedy. The question is whether the regulatory means chosen are reasonable in light of the nonpunitive objective.''
\end{quote}
\noindent --- Justice Kennedy, for the Court, 538 U.S. at 105.

\begin{quote}
``The process is more analogous to a visit to an official archive of criminal records than it is to a scheme forcing an offender to appear in public with some visible badge of past criminality. The Internet makes the document search more efficient, cost effective, and convenient for Alaska's citizenry.''
\end{quote}
\noindent --- Justice Kennedy, for the Court, 538 U.S. at 99.

\end{document}
